% Native author continuation, 2026-10-20. No inherited proof is modified.
\section{Two-sided capacity of the actual marked angular source}
\label{sec:mb-capacity}

The source in this section is the original one-copy marked flight, not a
replacement collision process.  A circle of directions is attached to
each of the two sides of the marked disk.  The disk is translated to
the origin only to name its line coordinates.  Following the line
backwards through $j+1$ true collisions fixes the original initial
disk; translating that disk to the fundamental cell recovers precisely
one original marked copy.  Following it forwards through $m-j$
collisions fixes the rest of the same record.  In particular the two
halves are not sampled independently.

Fix $R\in[R_-,R_+]$, $m\ge1$, $0\le j<m$, and
$u_0\in[0,3/200]$.  At $u_0=0$ use the original positive-side
unvisited trace: no collision is inserted at the marked foot.  Every
later first hit, including a later hit of the marked disk, remains
a collision.  At a positive level either prescribed incident side
may be used; if two sides are counted their charges are added.
The construction first uses all regular angular records.  Original
section, return, displacement, occupation and roof tests can
subsequently be retained as selectors on this unchanged source.

Keep the exact constants of Sections~\ref{sec:ag-source},
\ref{sec:uc-growth} and \ref{sec:nm-kernel}, and set
\begin{equation}\label{eq:mb-constants}
 \sigma=\frac{47}{59},\qquad
 E_*=\frac1{k_-a_\circ}=\frac{94}{\sqrt{1827}},\qquad
 C_{\rm foot}=\frac{28/5+R_+}{2R_-}=\frac{607}{90},
 \qquad Q_{\rm mb}=C_{\rm foot}H_{\rm nm}E_* .
\end{equation}
The enclosing horizon $28/5$ is a proved upper bound for the
original physical flights.  None of these constants depends on a
word, collision count, clearance, radius, or exact record.

\subsection{Suprema on genuine one-sided source intervals}

For $l\ge1$, let $P_l^\pm$ be the original momenta at the $l$th
forward or backward collision from the foot, with the sign convention
of Theorem~\ref{thm:nm-kernel}.  On a regular interval put
$A_l^\pm=|\partial_\vartheta P_l^\pm|$.
The following strengthens the pointwise root bound: it sums
suprema on a complete physical partition before choosing any target.

\begin{theorem}[One-sided physical branch capacity]
\label{thm:mb-branch-sup}
For either time direction there is a countable partition
$\mathcal P_l^\pm$ of the two angular circles, up to angular-null
sets, into disjoint open source intervals.  Each interval has a
fixed actual lifted half-word.  Its first $l-1$ collisions have
the homogeneous and length cuts used in the native growth theorem;
no terminal homogeneous subdivision is made.  With
\[
 a_I^\pm=\sup_{\vartheta\in I}(A_l^\pm(\vartheta))^{-1},
\]
one has
\begin{equation}\label{eq:mb-one-sided}
 \sum_{I\in\mathcal P_l^\pm}a_I^\pm\le H_{\rm nm},
 \qquad
 a_I^\pm\le E_*\sigma^{l-1}.
\end{equation}
These are suprema on actual intervals, not essential suprema of a
chosen version of the terminal momentum density.  Arbitrarily short
intervals are included.  The conclusion holds also for the specified
zero-clearance trace.
\end{theorem}
\begin{proof}
Choose the common angular seam before constructing either time
direction.  A translated seam in the reversal
$(\vartheta,\varsigma)\mapsto(\vartheta+\pi,-\varsigma)$ costs
exactly the same one seam in Lemma~\ref{lem:ag-branches}; it does
not introduce a second subdivision after the estimates.

For $l=1$, the first-hit and coordinate-seam intervals in that
lemma are the desired partition.  There are at most $M_*$
intervals for each side, and the exact first-hit derivative is
$A_1=\kappa a+c\ge k_-a_\circ$.  Hence the sum is at most
$2M_*E_*\le H_{\rm nm}$.

For $l\ge2$, use the source-labelled first-image family
$\rho_C\,d\ell_C$ of Theorem~\ref{thm:ag-family}.
Its density suprema sum to at most $A_*$.  For each $C$ apply
Theorem~\ref{thm:uc-growth} through exactly $l-2$ further real
collisions.  Let $V$ be one resulting short homogeneous curve and
$J_V(T_R^{-(l-2)})$ its inverse arclength Jacobian.  Cut $V$
only at the finitely many next-hit branches required for the last
actual collision.  There are at most $J_f$ of them.  On each such
branch the exact terminal cancellation gives
$|dp_{\rm next}/d\ell_V|\ge k_-t_*$.
The inverse angular derivative on its original source interval is
therefore bounded by
\[
 \frac{\|\rho_C\|_\infty}{k_-t_*}
                   \sup_V J_V(T_R^{-(l-2)}).
\]
Sum this bound over the last-hit cuts, the complete growth partition,
and the labelled first-image curves.  The result is
$A_*C_gJ_f/(k_-t_*)\le H_{\rm nm}$.
Every interval here is a pullback along injective regular curve
branches.  Homogeneity and length cuts preserve connectedness on
those graphs.  Thus these are disjoint intervals in the original
angular source, not a multiset of alternative representations.
The countable singular-null coverage is exactly the one proved
in Theorem~\ref{thm:uc-growth}.  No derivatives or component counts
of later selectors occur.

For the second inequality use the actual positive optical product
in \eqref{eq:bf-half-product}.  Its first free-flight matrix has
$B=a\ge a_\circ$ and $D=1$.  Each additional interior collision
and following flight replaces these entries by
\[
 B_{\rm new}=(1+\ell k)B+\ell D,\qquad
 D_{\rm new}=kB+D,\qquad
 k=\frac{2}{Rc}>0.
\]
Here $\ell\ge t_*=3/50$ and $c\le1$, so
$1+\ell k\ge1+2t_*/R_+=59/47$.  Positivity gives
$B\ge a_\circ(59/47)^{l-1}$.
The endpoint is incoming and has no extra kick in this product:
$A_l=\kappa B+c_lD\ge k_-B$.
This proves the pointwise bound and hence its interval supremum.
It remains valid as any positive interior incidence decreases;
no upper bound for a grazing kick was used.

At $u_0=0$ perform the same construction directly on the strict
positive-side source of Theorem~\ref{thm:nm-trace}.  The foot is
outside every other disk by the strict standoff estimate.  The
first-hit intervals, their exact derivatives, the factor $c$ in
$\rho_C$, and all curve-class estimates of 279--283 hold at this
level as well.  Later genuine tangencies exclude only the countable
angular set in the finite-time growth proof.  The skipped tangency
at the designated foot is fixed throughout and is never used as a
reflection.  This constructs the partition and proves its suprema
at the trace itself; a weak convergence assertion is not used to
transfer a sum of suprema.  Reversal is the same arclength and
angular isometry as in Theorem~\ref{thm:nm-kernel}, with the original
selectors transported through it.
\end{proof}

\subsection{An interval overlay estimate without a word-count factor}

\begin{lemma}[Weighted intersections of two interval families]
\label{lem:mb-overlay}
Let $\mathcal I,\mathcal J$ be two countable families of disjoint
open subintervals of one fixed bounded interval.  For nonnegative
weights $a_I,b_J$, with finite sums $A=\sum_Ia_I$, $B=\sum_Jb_J$,
put $a_*=\sup_Ia_I$ and $b_*=\sup_Jb_J$, using zero for an
empty family.  Then
\begin{equation}\label{eq:mb-overlay}
 \sum_{\substack{I\in\mathcal I,\ J\in\mathcal J\\I\cap J\ne\varnothing}}
                  a_Ib_J\le b_*A+a_*B .
\end{equation}
The same assertion holds on disjoint labelled copies of the interval,
using the combined sums and suprema.
\end{lemma}
\begin{proof}
Write an interval's left endpoint as $\ell_I$ or $\ell_J$.
For each nonempty intersection charge its product to $I$ if
$\ell_I\ge\ell_J$, and to $J$ otherwise.  A fixed $I$ receives
at most one charge: two disjoint open $J$ intervals starting no
later than $\ell_I$ cannot both intersect $I$ immediately to the
right of $\ell_I$.  The same argument applies to $J$ when its
left endpoint is the larger one.  Equal endpoints are assigned
to $I$ only.  Thus charges to $I$ sum to at most $b_*A$, and
charges to $J$ to at most $a_*B$.  Apply this first to finite
subfamilies.  Every finite collection of intersections is contained
in such subfamilies; nonnegative exhaustion proves the countable
assertion.  Summing over labelled components proves the last claim.
No lower interval length or connectedness of a union of intervals
is required.
\end{proof}

For the original $j$th marked flight overlay
$\mathcal P_{j+1}^-$ and $\mathcal P_{m-j}^+$ on each labelled
angular circle.  Denote the nonempty intersections by
$\mathcal O_{m,j,R,u_0}$.  An element has one complete genuine
collision word and its indicated interior homogeneous cuts.
Different intervals can carry the same word.  They are not merged
by assuming that its real physical domain is connected.

\begin{theorem}[Complete two-sided marked inverse-weight capacity]
\label{thm:mb-two-sided-sum}
Let $u_x=|\partial_{\alpha_0}p_m|^{-1}$ be the original inverse
endpoint-momentum weight on the complete regular marked source.
Then
\begin{equation}\label{eq:mb-two-sided-sum}
 \sum_{O\in\mathcal O_{m,j,R,u_0}}\sup_O u_x
       \le Q_{\rm mb}\bigl(\sigma^j+\sigma^{m-j-1}\bigr).
\end{equation}
This includes both mark signs and all actual half-words and their
intersections, also at the one-sided zero-clearance level.
Any original endpoint cap or exact-record selector can be retained
inside the positive sum.  Selectors are restricted within the
existing intervals and do not create extra charges.
Consequently, for arbitrary specified levels $u_j\in[0,3/200]$,
\begin{equation}\label{eq:mb-depth-sum}
 \sum_{j:\,\min(j,m-1-j)\ge D}
 \ \sum_{O\in\mathcal O_{m,j,R,u_j}}\sup_O u_x
       \le\frac{2Q_{\rm mb}}{1-\sigma}\sigma^D .
\end{equation}
An empty index set contributes zero.
\end{theorem}
\begin{proof}
On every actual overlay interval, Theorem~\ref{thm:bf-foot-determinant}
and Corollary~\ref{cor:bf-wronskian} give
\[
 u_x=\frac1{R A_-^pA_+^p(h_-+h_+)},\qquad
 h_-+h_+\ge\frac2{\tau_++R_+}.
\]
The two time directions fix the opposite clearance-column signs
in this determinant.  They yield a positive sum, not a difference
that could vanish.  Hence on $O=I\cap J$,
$\sup_Ou_x\le C_{\rm foot}a_I^-a_J^+$.
Apply Lemma~\ref{lem:mb-overlay} and
Theorem~\ref{thm:mb-branch-sup}.  The two sums of weights are at
most $H_{\rm nm}$; their respective maxima are at most
$E_*\sigma^j$ and $E_*\sigma^{m-j-1}$.
This gives \eqref{eq:mb-two-sided-sum} with exactly the constant
in \eqref{eq:mb-constants}.

The physical uniqueness in the one-mark flux identifies each
angular parameter and side with one actual marked copy at the
fixed index $j$.  Translating its starting disk does not add a
lattice multiplicity.  The overlay therefore sums the actual
source before imposing any record, roof or section selector.
Removing a test is used only to enlarge this positive upper bound.
Finally sum
$\sigma^j+\sigma^{m-j-1}$ over $D\le j\le m-1-D$.
Each of the two geometric sums is at most
$\sigma^D/(1-\sigma)$, proving \eqref{eq:mb-depth-sum}.
\end{proof}

\paragraph{Scope of the capacity statement.}
Equation~\eqref{eq:mb-two-sided-sum} is a genuine sum of physical
suprema on a complete marked angular source.  It does not follow
by choosing one common endpoint momentum for all intervals.
It supplies the comparison needed for the boundary theorem below,
including short and disconnected real word domains.
It is not the two-dimensional unrestricted pressure
$\sum_w\sup_{[w]}j_m$ in
Input~\ref{input:v112-geometric-pressure}, and is not yet a
pointwise roof density.  The original physical roof, its interval
widths, all four recorded coordinates and every outer limit remain
unchanged.
